\section{Main Results}\label{sec main}
\cmt{\begin{definition}[Linear Gaussian Problem (LGP)] \label{def LGP}Given latent vector $\bt^\st\in\R^d$, covariance $\bSi$ and noise level $\sigma$, suppose each example in $\Sc$ is generated independently as $y_i=\x_i^T\bt^\st+ z_i$ where $z_i\sim\Nn(0,\sigma^2)$ and $\x_i\sim\Nn(0,\bSi)$. Additionally, the map $\phi(\cdot)$ is identity and $p=d$.
\end{definition}}

\cmt{Given an ERM problem \eqref{erm} with nonlinear map $\phi$, we will relate it to an equivalent LGP to characterize its properties. The equivalence is ensured by setting up LGP to exhibit similar second order characteristics as the original problem.
\begin{definition}[Equivalent Linear Problem] Given distribution $(\x,y)$ $\sim\Dc$, the equivalent LGP($\bt,\bSi,\sigma$) with $n$ samples is given with parameters $\bSi=\E[\x\x^T]$, $\bt^\st=\bSi^{-1}\E[y\x]$ and $\sigma=\E[(y-\x^T\bt^\st)^2]^{1/2}$. 
\end{definition}
%from equivalent LGP to the original problem  %It should be emphasized that it is important
Our key theoretical contribution is formalizing the DC of the LGPs which also enables us to characterize pruning/retraining dynamics. We then empirically verify that DC and pruning dynamics of equivalent LGPs can successfully predict the original problem which can be random feature regression. In the context of high-dimensional learning literature, equivalent LGP is known to be a useful proxy for predicting the performance of the original ERM task \cite{}. This work goes beyond the existing works by demonstrating it can also successfully predict pruning/retraining dynamics. Formalizing this connection is an important future research avenue and it can presumably be accomplished by building on the recent high-dimensional universality results such as \cite{}.}

\cmt{
\noindent\textbf{Covariance diagonalization:} Let $\bSi$ have eigenvalue decomposition $\Ub\bar{\bSi}\Ub^T$. Let $\bth$ be the minimizer of empirical risk. Observe that $\X=\bar{\X}\Ub^T$ where $\bar{\X}$ has features with diagonal covariance $\bar{\bSi}$. Let $\btbh$ be the solution of ERM for $(\y,\bar{\X})$ pair. Then $\bth=\Ub\btbh$.
% subject to sparsity constraint i.e.~solving


\begin{definition}[Linear model] \label{d model}Fix a positive semidefinite covariance matrix $\bSi\in\R^{p\times p}$ and latent parameter $\bt\in\R^p$ satisfying $\bt^T\bSi\bt=\alpha$. Suppose $(\x_i)_{i=1}^n\distas\Nn(0,\bSi)$. The samples $(\x_i,y_i)_{i=1}^n$ are generated as 
\[
y_i=\x_i^T\bt+\sigma z_i,
\]
where $\sigma^2$ is the noise variance and $(z_i)_{i=1}^n$ is the noise sequence.
\end{definition}

Declare the first $s$ entries of $\bt$ to be $\tb=\Fb_s(\bt)$. We estimate $\tb$ via
\begin{align}
\tbh=\Fb_s(\bth)\quad\text{where}\quad \bth=\arg\min_{\bt'} \tn{\y-\X\bt'}.\label{optim me}
\end{align}
If the problem is over-parameterized ($p>n$), there are infinitely many feasible solutions $\bth$ achieving zero empirical loss. In this case, we investigate the minimum $\ell_2$ norm solution which is pseudo-inverse given by
\begin{align}
\tbh=\Fb_s(\bth)\quad\text{where}\quad \bth=\arg\min_{\bt'} \tn{\bt'}\quad\text{subject to}\quad \y=\X\bt'.\label{optim me2}
\end{align}
For the following discussion we set $\X=\Xb\sSi$ so that $\Xb\distas\Nn(0,1)$.}

%%%%%%%%%%%%%%%%%%%%%%%%%%%%%%%%%%%%%%%%%%%%%

Here, we present our main theoretical result: a sharp asymptotic characterization of the distribution of the solution to overparameterized least-squares for correlated designs. We further show how this leads to a sharp prediction of the risk of magnitude-based pruning.  
Concretely, for the rest of this section, we assume the linear Gaussian problem (LGP) of Definition \ref{def LGP}, the overparameterized regime $k=p>n$ and the min-norm model
%\begin{align}
%\hat{\betab} = \left\{ \arg\min_{\betab} \tn{\betab} ~\text{s.t.}~\tn{\y-\X\betab}=\min_{\betab'}\tn{\y-\X\betab'} \right\}\label{eq:min_norm}
%\end{align} 
\begin{align}
\hat{\betab} = \arg\min_{\betab} \tn{\betab} ~\text{s.t.}~\y=\X\bt.\label{eq:min_norm}
\end{align} 
As mentioned in Sec.~\ref{sec:setup}, $\hat{\betab}$ is actually given in closed-form as $\bth=\X^\dagger \y$. Interestingly, our analysis of the distribution of $\bth$ does not rely on the closed-form expression, but rather follows by viewing $\bth$ as the solution to the convex linearly-constrained quadratic program in \eqref{eq:min_norm}. Specifically, our analysis uses the framework of the convex Gaussian min-max Theorem (CGMT) \cite{thrampoulidis2015regularized}, which allows to study rather general inference optimization problems such as the one in \eqref{eq:min_norm}, by relating them with an auxiliary optimization that is simpler to analyze \cite{StoLASSO,oymak2013squared,thrampoulidis2015regularized,thrampoulidis2018precise,salehi2019impact,taheri2020fundamental}.
%
Due to space considerations, we focus here on the more challenging overparameterized regime and defer the analysis of the underparameterized regime to the SM. 

%Note that $\hat{\betab}$ is the LS solution when $p<n$ and the min-norm interpolator when $p>n$. \so{Due to space considerations, here, we focus on the latter, more challenging, overparameterized regime and defer the underparameterized regime to SM.}
%%Here, we focus on the latter, more challenging, overparameterized regime.


\subsection{Distributional Characterization of the Overparameterized Linear Gaussian Models}

\noindent\emph{Notation:}  We first introduce additional notation necessary to state our theoretical results. $\odot$ denotes the entrywise product of two vectors and $\onebb_p$ is the all ones vector in $\R^p$.
%
 The \emph{empirical distribution} of a vector $\x\in\R^p$ is given by 
$ \frac{1}{p}\sum_{i=1}^p \delta_{x_i}$, where $\delta_{x_i}$ denotes a Dirac delta mass on $x_i$. Similarly, the empirical joint distribution of vectors $\x, \x'\in \R^p$ is $\frac{1}{p} \sum_{i=1}^p \delta_{(x_i, x'_i)}$. The \emph{Wasserstein-$k$} ($W_k$) distance between two measures $\mu$ and $\nu$ is defined as 
%Given two probability measures $\mu$ (on a space $\mathcal{X}$) and $\nu$ (on a space $\mathcal{Y}$), a coupling $\rho$ of $\mu$ and $\nu$ is a probability distribution on $\mathcal{X}\times\mathcal{Y}$ whose marginals coincide with $\mu$ and $\nu$, respectively. 
%For $k\ge 1$, the \emph{Wasserstein-$2$} ($W_2$) distance between two probability measures $\mu$, $\nu$ on $\R^n$ is defined by
%
%\begin{align}
%
$W_k(\mu,\nu) \equiv \left(\inf_{\rho}\E_{(X,Y)\sim \rho}|X-Y|^k\right)^{1/k},$ %\label{eq:WassersteinDef}
%
%\end{align}
%
where the infimum is over all the couplings of $\mu$ and $\nu$, i.e. all random variables $(X,Y)$ such that $X\sim\mu$ and $Y\sim\nu$ marginally.
%X ? µ and Y ? ? marginally
A sequence of probability distributions $\nu_p$ on $\R^m$ {converges in $W_k$} to 
$\nu$, written $\nu_p\stackrel{W_k}{\Longrightarrow} \nu$, if $W_k(\nu_p,\nu) \rightarrow 0$ as $p \rightarrow \infty$. Finally, we say that a function $f:\R^m\rightarrow\R$ is pseudo-Lipschitz of order $k$, denoted $f\in\rm{PL}(k)$, if there is a constant $L>0$ such that for all $\x,\y\in\R^m$,
$
|f(\x) - f(\y)|\leq L(1+\tn{\x}^{k-1}+\tn{\y}^{k-1})\|\x-\y\|_2.
$
\fy{We call $L$ the $\rm{PL}(k)$ constant of $f$.}
%\somm{Is this convergence in probability or almost sure?} \CT{I understand that here convergence is a.s.. All our results are in probability / wpa 1. So formally, speaking, we do not show $W_k$ convergence. But is ok because the theorem is stated in terms of convergence of test functions.} 
An equivalent definition of $W_k$ convergence is that, for any $f\in\rm{PL}(k)$, $\lim_{p\rightarrow\infty}\E f(X_p)=\E f(X)$, where expectation is with respect to $X_p\sim\nu_p$ and $X\sim\nu$. For a sequence of random variables $\mathcal{X}_{p}$ that converge in probability to some constant $c$ in the limit of Assumption \ref{ass:linear} below, we  write $\mathcal{X}_{p}\rP c$. 

%We say that a sequence of probability distributions $\nu_p$ on $\R^m$ {converges in $W_k$} to $\nu$, written $\nu_p\stackrel{W_k}{\Longrightarrow} \nu$, if $W_k(\nu_p,\nu) \rightarrow 0$ as $p \rightarrow \infty$. An equivalent definition is that, for any $f\in\rm{PL}(k)$, $\lim_{p\rightarrow\infty}\E f(X_p)=\E f(X)$, where expectation is with respect to $X_p\sim\nu_p$ and $X\sim\nu$ (e.g., \cite{montanari2017estimation}).

\vp
Next, we formalize the set of assumption under which our analysis applies. Our asymptotic results hold in the linear asymptotic regime specified below.
\begin{assumption}\label{ass:linear} We focus on a double asymptotic regime where $n,p,s\rightarrow\infty$ at fixed overparameterization ratio $\kappa:=p/n>0$ and sparsity level $\alpha:=s/p\in(0,1)$.  
\end{assumption}

%
Additionally, we require certain mild assumptions on the behavior of the covariance matrix $\Sigmab$ and of the true latent vector $\bt^\star$. For simplicity, we assume here that $\Sigmab = \diag{[\Sigmab_{1,1},\ldots,\Sigmab_{p,p}]}$.% and we also need that $\Sigmab_{i,i}$ are non-zero and bounded. 
%Moreover, we specify the following two assumptions.

\begin{restatable}{assumption}{asstwo}\label{ass:inv} {The covariance matrix $\Sigmab$ is diagonal} and there exist constants $\Sigma_{\min},\Sigma_{\max}\in(0,\infty)$ such that:
$
\Sigma_{\min}\leq\bSi_{i,i}\leq \Sigma_{\max},
$
for all $i\in[p].$
\end{restatable}


%Finally, we assume a limit on the joint empirical measure of $(\diag{\Sigma},\betas)$ as follows.
\begin{restatable}{assumption}{assthree}\label{ass:mu}
% Let $\nub_i:= {{p}(\bt^\star_i)^2{\lai}},~i\in[p]$. The joint empirical distribution of $\{(\lai,\nub_i)\}_{i\in[p]}$ converges in Wasserstein-2 distance to a probability distribution $\mu$ on $\R_{>0}\times\R_{>0}$, i.e.,
%$
%\frac{1}{p}\sum_{i\in[p]}\delta_{(\lai,\nub_i)} \stackrel{W_2}{\Longrightarrow} \mu.
%$
The joint empirical distribution of $\{(\lai,\sqrt{p}\betas_i)\}_{i\in[p]}$ converges in {Wasserstein-k} distance to a probability distribution $\mu$ on $\R_{>0}\times\R$ {for some {$k\geq 4$}}. That is
$
\frac{1}{p}\sum_{i\in[p]}\delta_{(\lai,\sqrt{p}\betas_i)} \stackrel{W_k}{\Longrightarrow} \mu.
$
\end{restatable}

With these, we are ready to define, what will turn out to be, the asymptotic DC in the overparameterized regime.
\begin{definition}[Asymptotic DC -- Overparameterized regime]\label{def:Xi}
Let random variables $(\Lambda,B)\sim \mu$ (where $\mu$ is defined in Assumption \ref{ass:mu}) and fix $\kappa>1$. Define parameter $\xi$ as the unique positive solution to the following equation
\begin{align}\label{eq:ksi}
\E_{\mu}\Big[ \big({1+(\xi\cdot\Lambda)^{-1}}\big)^{-1} \Big] = {\kappa^{-1}}\,.
\end{align}
Further define the positive parameter $\gamma$ as follows:
\begin{align}
\label{eq:gamma}
\hspace{-0.1in}\gamma := 
%\frac{\sigma^2 + \E_{\mu}\left[\frac{N^2}{(\Lambda^{-1}+\xi)^2}\right]}{1-\kappa\E_{\mu}\left[\frac{1}{\left(1+(\xi\Lambda)^{-1}\right)^2}\right]}  
\Big({\sigma^2 + \E_{\mu}\Big[\frac{B^2\Lambda}{(1+\xi\Lambda)^2}\Big]}\Big)\Big/\Big({1-\kappa\E_{\mu}\Big[\frac{1}{\left(1+(\xi\Lambda)^{-1}\right)^2}\Big]}\Big).
\end{align}
%\ct{@Samet: In your notation $\gamma$ should be $\Gamma\kappa$. Everything matches except the second term in the numerator. In your notation, that second term becomes $\int_0^1\bz^2(t)\bt^2(t) \Sigmab(t) dt$ (note the extra $\Sigmab(t)$ compared to your formula)
%}
With these and $H\sim\Nn(0,1)$, define the random variable
\begin{align}
X_{\kappa,\sigma^2}(\Lambda,B,H) := \Big(1-\frac{1}{1+ \xi\Lambda}\Big) B + \sqrt{\kappa}\frac{\sqrt{\gamma}\,\Lambda^{-1/2}}{1+(\xi\Lambda)^{-1}} H, \label{eq:X}
\end{align}
and let $\Pi_{\kappa,\sigma^2}$ be its distribution.
\end{definition}

%\begin{align}
%&\Xi>0\quad\text{is the unique solution of}\quad 1=\int_{0}^1\frac{\kappa}{1+(\Xi\bSi(t))^{-1}}dt,\label{aux solve}\\
%&\Gamma=\frac{\sigma^2+\int_0^1\bz^2(t)\bt^2(t)dt}{\kappa(1-\int_0^1\frac{\kappa}{(1+(\Xi\bSi(t))^{-1})^2}dt)},\nn\\
%&\bz(t)=\frac{1}{1+\Xi\bSi(t)}\quad\text{and}\quad \bg(t)=\frac{\kappa\sqrt{\Gamma}}{1+(\Xi\bSi(t))^{-1}}\quad\text{for}\quad 0\leq t\leq 1.\nn
%\end{align}
%

%\begin{definition}\label{def:measure}
%Let random variables $(\Lambda,N)\sim \mu$ and $H\sim\Nn(0,1)$. Further let $\xi,\gamma$ as defined in Definition \ref{def:Xi}. 
%With these, define the random variable
%$$
%X_{\kappa,\sigma^2}(\Lambda,N,H) = \frac{-N}{\Lambda^{-1} + \xi} + \sqrt{\kappa}\frac{\sqrt\gamma\cdot\xi}{\Lambda^{-1} + \xi} H.
%$$
%Let $\Pi_{\kappa,\sigma^2}$ be the distribution of $X_{\kappa,\sigma^2}(\Lambda,N,H)$.
%\end{definition}


%\begin{definition}[Pseudo-Lipschitz test function] 
%We say that a function $f:\R^m\rightarrow\R$ is pseudo-Lipschitz of order $2$, denoted $f\in\rm{PL}(2)$, if there is a constant $L>0$ such that for all $\x,\y\in\R^m$,
%$$
%|f(\x) - f(\y)|\leq L\left(1+\tn{x}+\tn{y}\right)\|\x-\y\|_2.
%$$

%\end{definition}



% characterizing the DC of the ERM minimizer $\hat\bt$ in the overparameterized regime. 
%\so{Specifically, the theorem characterizes the distribution of the error vector weighted by the covariance matrix, i.e., of $\hat\w=\sqrt{\Sigmab}(\hat{\bt}-\bt^\star)$. We will show after the statement of the theorem that $\hat\w$ is exactly what matters for the characterization of the population loss of the pruned solution.}\som{not correct statement}
%Let $\cal{H}$ be the hard-thresholding operation. Draw $\hat{\bt}\sim \Dca(\bt,\bSi)$.
%\[
%\text{risk}_{\bSi^p}(\Hc(\hat{\bt}^p))\rightarrow \text{risk}_{\bSi}(\Hc(\hat{\bt}))
%\]
%\end{theorem}
%
%\begin{corollary} 
%Specifically, letting $\cal{H}:\R\rightarrow\R$ denote the hard-thresholding operation, it holds that
%\begin{align}
%\| \cal{H}( \hat\w^p )\|_2^2 \rP %\E\left[\left(\cal{H}\left(X_{\kappa,\sigma^2}(\Lambda,N,H) \right)\right)^2\right]
%\end{align}
%%\end{corollary}


\fy{Our main result establishes asymptotic convergence of the empirical distribution of $(\sqrt{p}\bth,\sqrt{p}\betas,\Sigmab)$ for a rich class of test functions. These are the functions within $\rm{PL}(3)$ that become $\rm{PL}(2)$ when restricted to the first two indices. Formally, we define this class of functions as follows
\begin{align}\label{eq:pdef}
\Plt :=\{&f:\R^2\times \Zc\rightarrow\R,~f\in \rm{PL}(3)~\text{and}~\\
&\sup_{z\in\Zc}\text{``}\rm{PL}(2)~\text{constant~of}~f(\cdot,\cdot,z)\text{''}<\infty\}.\nn
\end{align}
For pruning analysis, we set $\Zc=[\Sigma_{\min},\Sigma_{\max}]$ and define %the following functions%class are critical %has the form%It applies to all functions in $\rm{PL}(2)$ as well as the following class of functions
\begin{align}\label{eq:fdef}
\Fl:= \{ f:\R^2\times \Zc\rightarrow\R~\big|~&f(x,y,z) = z(y-g(x))^2\nn\\
&~\text{where}~g(\cdot)~\text{is Lipschitz}\}.
%f_\Lc(x,y,z) := \fx{z(y-g(x))}^2. %:\R\rightarrow \R
\end{align}
%where $g$ is Lipschitz. 
%\CT{I rewrote this as a family. Alternatively, Thm. 1 could be stated to hold for family of functions $f:\R^3\rightarrow\R$ such that $f\in\rm{PL}(3)$ and $g(x,y):=f(x,y,z)$ is $\rm{PL}(2)$ with some $\rm{PL}$-constant $L$ that uniformly over $z$ in the range of $f$.}
As discussed below, $\Fl$ is important for predicting the risk of the (pruned) model. In the SM, we prove that $\Fl\subset \Plt$. We are now ready to state our main theoretical result.}% which applies to $\Plt$
%\begin{theorem}[Asymptotic DC for LGP]\label{thm:master_W2} %Fix $\kappa=p/n>1$.
%Let Assumption \ref{ass:linear} hold with $\kappa>1$ and  further let Assumptions \ref{ass:inv} and \ref{ass:mu} hold. Consider $\hat{\bt}$ as in \eqref{eq:min_norm} and $\hat\Pi_n(\y,\X,\betas,\Sigmab):=\frac{1}{p}\sum_{i=1}^{p}\delta_{\sqrt{p}\hat{\bt}_i,\sqrt{p}\betas_i,\Sigmab_{i,i}}$, the joint empirical distribution of $(\sqrt{p}\hat\betab,\sqrt{p}\betas,\Sigmab)$. Recall the definition of the measure $\Pi_{\kappa,\sigma^2}$ in Definition \ref{def:Xi}. Then, $\hat\Pi_n(\y,\X,\betas,\Sigmab)$ converges in Wasserstein-k distance to $\Pi_{\kappa,\sigma^2}\otimes\mu$.
%Specifically, for any function $f:\R^3\rightarrow\R$, $f\in\rm{PL}(k)$ with $k\geq 3$, it holds that
%\begin{align}\label{eq:thm}
%\hspace{-0.1in}p^{-1} \sum_{i=1}^{p} f(\sqrt{p}\hat\betab_i,\sqrt{p}\betas_i,\Sigmab_{i,i}) \rP \E\left[f(X_{\kappa,\sigma^2},B,\Lambda) \right],
%\end{align}
%where the expectation is over $(\Lambda,B,H)\sim\mu\otimes\Nn(0,1).$
%\end{theorem}
% be any of the following two: (a) $f\in\rm{PL}(2)$, or, (b) $f=f_\Lc$ defined in \eqref{eq:fdef} \fx{where $g$ is a Lipschitz function}. 
\cts{
\begin{restatable}[Asymptotic DC -- Overparameterized LGP]{theorem}{mainthm}\label{thm:master_W2} %Fix $\kappa=p/n>1$.
Fix $\kappa>1$ and suppose Assumptions \ref{ass:inv} and \ref{ass:mu} hold. Recall the solution $\bth$ from \eqref{eq:min_norm} and let 
\[
\hat\Pi_n(\y,\X,\betas,\Sigmab):=\frac{1}{p}\sum_{i=1}^{p}\delta_{(\sqrt{p}\bth_i,\sqrt{p}\betas_i,\Sigmab_{i,i})}
\] be the joint empirical distribution of $(\sqrt{p}\bth,\sqrt{p}\betas,\Sigmab)$. Let $f:\R^3\rightarrow\R$ be a function in $\Plt$ defined in \eqref{eq:pdef}. We have that
%\begin{align}\label{eq:thm_app_f}
%\frac{1}{p} \sum_{i=1}^{p} f(\sqrt{p}\bth_i,\sqrt{p}\betas_i,\Sigmab_{i,i}) \rP \E_{(\Lambda,B,H)\sim\mu\otimes\Nn(0,1)}\left[f(X_{\kappa,\sigma^2},B,\Lambda) \right].
%\end{align}
\begin{align}\label{eq:thm}
\hspace{-0.1in}\frac{1}{p} \sum_{i=1}^{p} f(\sqrt{p}\hat\betab_i,\sqrt{p}\betas_i,\Sigmab_{i,i}) \rP \E\left[f(X_{\kappa,\sigma^2},B,\Lambda) \right].
\end{align}
\end{restatable}
}


%\begin{restatable}[Goldbach's conjecture]{theorem}{goldbach}
%\label{thm:goldbach}
%Every even integer greater than 2 can be expressed as the sum of two primes.
%\end{restatable}
%
%\goldbach*

%%%%%%%%%%%%%%%%%%%%%%%%%%%%%%%%%%%%%%%%%%%%%
%Note that the integer $k$ appearing in the statement of the Theorem \ref{thm:master_W2} is the same 
%\begin{remark}(Novelty) Compare to Surprises. Compare to CGMT.
%\end{remark}

As advertised, Theorem \ref{thm:master_W2} fully characterizes the joint empirical distribution of the min-norm solution, the latent vector and the covariance spectrum. The asymptotic DC allows us to precisely characterize several quantities of interest, such as estimation error, generalization error etc.. For example, a direct application of \eqref{eq:thm} to the function $f(x,y,z)=z(y-x)^2\in \Fl\subset\Plt$ directly yields the risk prediction of the min-norm solution recovering \cite[Thm.~3]{hastie2019surprises} as a special case. Later in this section, we show how to use Theorem \ref{thm:master_W2} towards the more challenging task of precisely characterizing the risk of magnitude-based pruning.

Before that, let us quickly remark on the technical novelty of the theorem. Prior work has mostly applied the CGMT to isotropic features. Out of these, only very few obtain DC, \cite{thrampoulidis2018symbol,miolane2018distribution}, while the majority focuses on simpler metrics, such as squared-error. Instead, Theorem \ref{thm:master_W2} considers correlated designs and the overparameterized regime. The most closely related work in that respect is \cite{montanari2019generalization}, which very recently obtained the DC of the max-margin classifier. Similar to us, they use the CGMT, but their analysis of the auxiliary optimization is technically different to ours. Our approach is similar to \cite{thrampoulidis2018symbol}, but extra technical effort is needed to account for correlated designs and the overparameterized regime.

\subsection{From DC to Risk Characterization}\label{sec:risk}
%In the remaining of this section, we show how to use Theorem \ref{thm:master_W2} to precisely characterize the risk of the magnitude-based pruning. 
First, we consider a simpler ``threshold-based" pruning method that applies a fixed threshold at every entry of $\hat\betab$. Next, we relate this to magnitude-based pruning and obtain a characterization for the performance of the latter.
%
In order to define the threshold-based pruning vector, let \[
\Tc_t(x) = \begin{cases} x & \text{if } |x|>t \\ 0 & \text{otherwise} \end{cases},
\] be the hard-thresholding function with fixed threshold $t\in\R_+$. Define $\bth^\Tc_{t} := \Tc_{t/\sqrt{p}}(\hat\bt)$, where $\Tc_t$ acts component-wise. Then, the population risk of   $\bth^\Tc_t$ becomes
\begin{align}
\Lc(\bth^\Tc_t) &= \E_{\Dc}[(\x^T(\bt^\star-\bth^\Tc_t) + \sigma z)^2] \nn\\
 \cmt{&= \sigma^2 + (\bt^\star-\bth^\Tc_t)^T\Sigmab(\bt^\star-\bth^\Tc_t) \nn \\}
  &= \sigma^2 + \frac{1}{p}\sum_{i=1}^p\Sigmab_{i,i}\big(\sqrt{p}\betas_i-\Tc_{t}(\sqrt{p}\hat\betab_i)\big)^2\nn\\
  &\rP \sigma^2 + \E\left[ \Lambda\left(B-\Tc_t(X_{\kappa,\sigma^2})\right)^2 \right]\,.
  \label{eq:loss_w2}
%\\ &= \sigma^2 + \tn{\Tc_t\left(\sqrt{\Sigmab}(\hat\bt-\bt^\star)\right)}^2 + \tn{\sqrt{\Sigmab}(\bt^\star-\Tc_t(\bt^\star))}^2\nn
%\\ &= \sigma^2 + \frac{1}{p}\tn{\Tc_t(\sqrt{\p}\hat\w)}^2 + \tn{\sqrt{\Sigmab}(\bt^\star-\Tc_t(\bt^\star))}^2 \label{eq:loss_w2}
\end{align}
%In the first equality, we let $z\sim\Nn(0,1)$. 
%In the second, we used $\E[\x\x^T]=\Sigmab$. In the third line, we wrote $\bt^\star = \Tc_t(\bt^\star) + (\bt^\star - \Tc_t(\bt^\star))$ and we also used the fact that $\Sigmab$ is diagonal so that $\sqrt{\Sigmab}\Tc_t(\vb)=\Tc_t(\sqrt{\Sigmab}\vb)$ for any vector $\vb\in\R^p$. In the last line, we recalled the definition of $\hat\w$ in Theorem \ref{thm:master_W2}. 
In the second line above, we note that $\sqrt{p}\Tc_{t'}(x)=\Tc_{\sqrt{p}t'}(\sqrt{p}x)$. In the last line, we apply \eqref{eq:thm}, after recognizing that the function $(x,y,z)\mapsto z(y-\Tc_t(x))^2$ 
%is pseudo-Lipschitz of order $3$ (see SM \fx{Lemma \ref{lem:fL}} for a proof)%\footnote{Notice that the function $\Tc_t$ is discontinuous. However, it can be %approximated arbitrarily closely with Lipschitz functions.}. \fx{Additionally observe %that this function 
is a member of the $\Fl$ family defined in \eqref{eq:fdef}.
%} 
As in \eqref{eq:thm}, the expectation here is with respect to $(\Lambda,B,H)\sim\mu\otimes\Nn(0,1)$. 
%\somm{Here, more justification on why Lip approx is sufficient would be great. Obviously Lip approx should be fine because asymptotic distribution is continuous.}
%\so{$\E[\Tc_t(X_{\kappa,\sigma^2}(\Lambda,N,H))^2]$. 
%The third term converges to }\som{Forgot covariance multiplier?}

Now, we show how to use \eqref{eq:loss_w2} and Theorem \ref{thm:master_W2} to characterize the risk of the magnitude-based pruned vector $\betab^M_s:=\Tb_s(\bth)$. Recall, here from Assumption \ref{ass:linear} that $s=\alpha p$. To relate $\hat{\betab}^M_s$ to $\bth^\Tc_t$, consider the set  $\Sc_t:=\{i\in[p]\,:\,\sqrt{p}|\hat\betab_i| \geq t \}$ for some constant $t\in\R_+$ (not scaling with $n,p,s$). Note that the ratio ${|\Sc_t|}/{p}$ is equal to
\begin{align}
p^{-1}\sum_{i=1}^p\mathbb{1}_{[\sqrt{p}|\hat\betab_i| \geq t]} \rP \E[\mathbb{1}_{[|X_{\kappa,\sigma^2}|\geq t]}] = \P\left(|X_{\kappa,\sigma^2}|\geq t\right).\nn
\end{align}
Here, $\mathbb{1}$ denotes the indicator function and the convergence follows from Theorem \ref{thm:master_W2} when applied to a sequence of bounded Lipschitz functions approximating the indicator. Thus, by choosing 
\begin{align}\label{eq:tstar}
t^\star:=\sup\left\{t\in\R\,:\, \P(|X_{\kappa,\sigma^2}|\geq t) \geq \alpha \right\},
\end{align}
it holds that ${|\Sc_t|}/{p}\rP\alpha$. In words, and {observing that $X_{\kappa,\sigma^2}$ admits a continuous density (due to the Gaussian variable $H$)}: for any  $\eps>0$, in the limit of $n,p,s\rightarrow\infty$, the vector $\bth^\Tc_{t^\star}$ has $(1\pm\eps) \alpha p= (1\pm\eps) s$ non-zero entries, which correspond to the largest magnitude entries of $\bth$, with probability approaching1. Since this holds for arbitrarily small $\eps>0$, recalling $t^\star$ as in \eqref{eq:tstar}, we can conclude from \eqref{eq:loss_w2} that the risk of the magnitude-pruned model converges as follows.
{\begin{restatable}[Risk of Magnitude-pruning]{corollary}{cormag}\label{cor:mag}
Let the same assumptions and notation as in the statement of Theorem \ref{thm:master_W2} hold. Specifically, let $\hat\betab$ be the min-norm solution in \eqref{eq:min_norm} and $\bth_s^M:=\Tb_s(\bt)$ the magnitude-pruned model at sparsity $s$. Recall the threshold $t^\st$ from \eqref{eq:tstar}. The risk of $\bth^M_s$ satisfies the following in the limit of $n,p,s\rightarrow\infty$ at rates $\kappa:=p/n>1$ and $\alpha:=s/p\in(0,1)$ (cf. Assumption \ref{ass:linear}):
\begin{align}
\Lc(\bth^M_s) \rP \sigma^2 + \E\left[ \Lambda\left(B-\Tc_{t^\star}(X_{\kappa,\sigma^2})\right)^2 \right],\nn
\end{align}
where the expectation is over $(\Lambda,B,H)\sim\mu\otimes\Nn(0,1).$
\end{restatable}}
%\begin{align}
%\Lc(\bth^M_s) \rP \sigma^2 + \E\left[ \Lambda\left(B-\Tc_{t^\star}(X_{\kappa,\sigma^2})\right) \right].\nn
%\end{align}
%\somm{One way to fully rigorize this line and Lip approx is showing that any eps-fraction of entries of $||\hat{\beta}||$ is small in $\eps$. With that, it becomes crystal clear I think.}
%where .

\cmt{LET US NOT FORGET TO ADD OR COMMENT ABOUT UNDERPARAM!}
% Prediction of
\subsection{Non-asymptotic DC and Retraining Formula}
While Theorem \ref{thm:master_W2} is stated in the asymptotic regime, during analysis, the DC arises in a non-asymptotic fashion. The following definition is the non-asymptotic counterpart of Def.~\ref{def:Xi}. We remark that this definition applies to arbitrary covariance (not necessarily diagonal) by applying a simple eigen-rotation before and after the DC formula associated with the diagonalized covariance.
\cmt{
\begin{definition}[Non-asymptotic DC] \label{aux_def}Fix $p>n\geq 1$ and set $\kappa=p/n>1$. Given $\sigma>0$, covariance $\bSi=\Ub\text{diag}(\bla)\Ub^T$ and latent vector $\bt$, set $\bar{\bt}=\Ub^T\bt$ and define the unique non-negative terms $\Xi,\Gamma,\bz\in\R^p$ and $\bg\in\R^p$ as follows\vspace{-0pt}
\begin{align}
&\Xi>0\quad\text{is the solution of}\quad 1=\frac{\kappa}{p}\sum_{i=1}^p\frac{1}{1+(\Xi\bla_i)^{-1}},\nn\\
&\Gamma=\frac{\sigma^2+\sum_{i=1}^p\bla_i\bz_i^2\bar{\bt}_i^2}{\kappa(1-\frac{\kappa}{p}\sum_{i=1}^p{(1+(\Xi\bla_i)^{-1})^{-2}})},\nn\\
&\bz_i=\frac{1}{1+\Xi\bla_i}\quad\text{and}\quad \bg_i=\frac{\kappa\sqrt{\Gamma}}{1+(\Xi\bla_i)^{-1}}\quad\text{for}\quad 1\leq i\leq p.\nn
\end{align}
The non-asymptotic OP-LS distribution is defined as the following $\Ub$-rotated normal distribution
\[
\Dc_{\sigma,\bSi,\bt}=\Ub\Nn((\onebb_p-\bz)\odot\bar{\bt},p^{-1}\text{diag}(\bla^{-1}\odot\bg^2)).
\]
\end{definition}}
%\somm{Justify this definition in more detail}
\begin{definition}[Non-asymptotic DC] \label{aux_def}Fix $p>n\geq 1$ and set $\kappa=p/n>1$. Given $\sigma>0$, covariance $\bSi=\Ub\text{diag}(\bla)\Ub^T$ and latent vector $\bt$, set $\bar{\bt}=\Ub^T\bt$ and define the unique non-negative terms $\xi,\gamma,\bz\in\R^p$ and $\bg\in\R^p$ as follows:\vspace{-0pt}
\begin{align}
&\xi>0\quad\text{is the solution of}\quad \kappa^{-1}={p^{-1}}\sum_{i=1}^p\big({1+(\xi\bla_i)^{-1}}\big)^{-1},\nn\\
&\gamma=\frac{\sigma^2+\sum_{i=1}^p\bla_i\bz_i^2\bar{\bt}_i^2}{1-\frac{\kappa}{p}\sum_{i=1}^p{(1+(\xi\bla_i)^{-1})^{-2}}},\nn\\
&\bz_i=({1+\xi\bla_i})^{-1}\quad\text{,}\quad \bg_i={\kappa\gamma}(1+(\xi\bla_i)^{-1})^{-2},~ 1\leq i\leq p.\nn
\end{align}
The non-asymptotic distributional prediction is given by the following $\Ub$-rotated normal distribution
\[
\Dc_{\sigma,\bSi,\bt}=\Ub\Nn((\onebb_p-\bz)\odot\bar{\bt},p^{-1}\text{diag}(\bla^{-1}\odot\bg)).
\]
\end{definition}
{We remark that this definition is similar in  spirit to the concurrent/recent work \cite{li2020exploring}. However, unlike this work, here we prove the asymptotic correctness of the DC, we use it to rigorously predict the pruning performance and also extend this to retraining DC as discussed next.} 

%E[x_Ix_I^T]^(-1)E[xx^T]E[xx^T]^(-1)E[yx] = E[x_Ix_I^T]^(-1)E[yx] = E[x_Ix_I^T]^(-1)E[yx_I]
\noindent \textbf{Retraining DC.} As the next step, we would like to characterize the DC of the solution after retraining, i.e.,~$\bth^{RT}$. We carry out the retraining derivation (for magnitude pruning) as follows. Let $\Ic\subset[p]$ be the nonzero support of the pruned vector $\bth^M_s$. Re-solving \eqref{eq:ERM} restricted to the features over $\Ic$ corresponds to a linear problem with effective feature covariance $\bSi_\Ic$ with support of non-zeros restricted to $\Ic\times \Ic$. For this feature covariance, we can also calculate the effective noise level and global minima of the population risk $\bt^\st_\Ic$. The latter has the closed-form solution $\bt^\st_\Ic=\bSi_\Ic^\dagger\bSi\bt^\st$. The effective noise is given by accounting for the risk change due to the missing features via $\sigma_\Ic=(\sigma^2+{\bt^\st}^T\bSi\bt^\st-{\bt^\st_\Ic}^T\bSi_\Ic\bt^\st_\Ic)^{1/2}$. With these terms in place, fixing $\Ic$ and using Def.~\ref{aux_def}, the retraining prediction becomes $\Dc_{\sigma_{\Ic},\bSi_{\Ic},\bts_\Ic}$. This process is summarized below.
\begin{definition}[Retraining DC] \label{RT_def}Consider the setting of Def.~\ref{aux_def} with $\sigma,\bSi,\bt^\st$ and sparsity target $s$. The sample $\bth^{RT}$ from the retraining distribution $\Dc^{\text{RT,s}}_{\sigma,\bSi,\bt^\st}$ is constructed as follows. Sample $\bth\sim \Dc_{\sigma,\bSi,\bt^\st}$ and compute the set of the top-$s$ indices $\Ic=\Ic(\Tb_s(\bth))$. Given $\Ic$, obtain the effective covariance $\bSi_\Ic\in\R^{p\times p}$, population minima $\bt^\st_\Ic\in\R^p$, and the noise level $\sigma_{\Ic}>0$ as described above. Draw $\bth^{RT}\sim\Dc_{\sigma_{\Ic},\bSi_{\Ic},\bts_\Ic}$.

\cmt{Obtain $\bSi_\Ic\in\R^{p\times p}$ by restricting the nonzero support of $\bSi$ to $\Ic\times \Ic$. Set $\Ic$ restricted population minima $\bt^\st_\Ic=\bSi_\Ic^\dagger\bSi\bt^\st$ and set the new noise level $\sigma_{\Ic}=(\sigma^2+{\bt^\st}^T\bSi\bt^\st-{\bt^\st_\Ic}^T\bSi_\Ic\bt^\st_\Ic)^{1/2}$. Draw $\bth^{RT}\sim\Dc_{\sigma_{\Ic},\bSi_{\Ic},\bt_\Ic}$. }
\end{definition}
%\som{Clarify here!}
Observe that, the support $\Ic$ depends on the samples $\Sc$ via $\bth$. Thus, our retraining DC is actually derived for the scenario when the retraining phase uses a fresh set of $n$ samples to break the dependence between $\Ic,\Sc$ (which obtains $\bth^{RT}_{\text{fresh}}$). Despite this, we empirically observe that the retraining DC predicts the regular retraining (reusing $\Sc$) performance remarkably well and perfectly predicts $\bth^{RT}_{\text{fresh}}$ as discussed in Figs.~\ref{figRF} and \ref{figRF2}. Finally, we defer the formalization of the retraining analysis to a future work. This includes proving that $\bth^{RT}_{\text{fresh}}$ obeys Def.~\ref{RT_def} asymptotically as well as directly studying $\bth^{RT}$ by capturing the impact of the $\Ic,\Sc$ dependency.
 
%\so{Here, the core technical challenges are addressing the randomness in $\Ic$ by characterizing its asymptotic behavior and properly mapping this behavior to the asymptotic retraining distribution.}

 %Additionally, if we indeed use fresh samples, it indeed exactly predicts the retraining performance (see Fig. \ref{figRF2}). 


%showing that asymptotically all $\Ic$'s will lead to the same retraining distribution
%theshowing that asymptotically all $\Ic$'s will lead to the same retraining distribution.








